% s03-1-4-characters.tex —— 3.1.4 小节：IFLE 的主角（印刷页 74–77）

\subsection{IFLE 的主角}
\label{sec:ifle-characters}

本小节先回顾 \cite[2.1 节]{ABC+24a} 中把 Kazhdan–Lusztig 范畴定义为因子化范畴的构造。

Kac–Moody 模的因子化范畴 $\widehat{\mathfrak{g}}\text{-}\mathrm{mod}_\kappa$ 带有因子化群
$LG$ 在层级 $\kappa$ 处的（强）作用。该因子化范畴与作用的构造见
\cite[附录 B.14]{ABC+24a}。因此下述定义有意义。

\begin{definition}
\label{def:kl}
因子化范畴 $\KL(G)_\kappa\in \FactAlgCat$ 为
\[
\widehat{\mathfrak{g}}\text{-}\mathrm{mod}^{L^+G}_\kappa .
\]
\end{definition}

可沿 $\Ran_{x_0}\to \Ran$ 限制范畴层 $\widehat{\mathfrak{g}}\text{-}\mathrm{mod}_\kappa$，
得到 $\widehat{\mathfrak{g}}\text{-}\mathrm{mod}_{\kappa,\Ran_{x_0}}$。由于显然的理由，
它是因子化 $\widehat{\mathfrak{g}}\text{-}\mathrm{mod}_\kappa$-模范畴。同样由于显然的
理由，$LG^{\Ran_{x_0}}$ 作用在 $\widehat{\mathfrak{g}}\text{-}\mathrm{mod}_{\kappa,\Ran_{x_0}}$
上。于是可取因子化 $L^+G$-模群 $I^{\Ran}_{x_0}$ 的不变量，得到一个因子化
$\KL(G)_\kappa$-模范畴，记作 $\IKL(G)_\kappa$。

\begin{definition}
\label{def:ikl}
$\FactModCat_{x_0}$ 中的对 $(\KL(G)_\kappa,\IKL(G)_\kappa)$ 由因子化范畴 $\KL(G)_\kappa$
与因子化 $\KL(G)_\kappa$-模范畴
\[
\IKL(G)_\kappa
=\widehat{\mathfrak{g}}\text{-}\mathrm{mod}^{I^{\Ran_{x_0}}}_{\kappa,\Ran_{x_0}}
\]
构成。
\end{definition}
\srcpage{74}

本小节余下部分及全文只关心 $\kappa=\kappa_{\mathrm{crit}}$ 的情形。

设 $\mathcal{C}$ 为 $\Ran$ 上带层级 $\kappa_{\mathrm{crit}}$ 处 $LG$-作用的范畴层，
则有幺半作用
\[
\Sph_G\otimes \mathcal{C}^{L^+G}\longrightarrow \mathcal{C}^{L^+G}.
\]
再设 $\mathcal{D}$ 为 $\Ran_{x_0}$ 上带层级 $\kappa_{\mathrm{crit}}$ 处 $LG$-作用的
范畴层，则有幺半作用
\[
\Aff_G\otimes \mathcal{D}^I\longrightarrow \mathcal{D}^I .
\]
于是可以把对 $(\KL(G)_{\mathrm{crit}},\IKL(G)_{\mathrm{crit}})$ 看作带有幺半对
$(\Sph_G,\Aff_G)$ 的幺半作用。
\srcpage{75}

下面转到谱侧。存在 $\check G$-opers 的仿射 $D$-概形 $\Op_{\check G}$，其定义见
\cite[3.1.3 节]{ABC+24a}。于是水平弧空间与水平环路空间的构造允许我们定义正则 opers 的
因子化空间
\[
\Op^{\reg}_{\check G}=L^+_\nabla \Op_{\check G}
\]
与亚纯 opers 的因子化空间
\[
\Op^{\mer}_{\check G}=L_\nabla \Op_{\check G}.
\]
有图表
\[
\begin{tikzcd}[column sep=large,row sep=large]
\Op^{\reg}_{\check G} \arrow[r] \arrow[d] & \Op^{\mer}_{\check G} \arrow[d]\\
\LS^{\reg}_{\check G} \arrow[r] & \LS^{\mer}_{\check G},
\end{tikzcd}
\]
它允许我们作出\emph{无单值 opers} 的因子化空间
\[
\Op^{\monfree}_{\check G}
=\Op^{\mer}_{\check G}\times_{\LS^{\mer}_{\check G}}\LS^{\reg}_{\check G}.
\]

\begin{definition}
\label{def:indcoh-opers}
无单值 opers 上归纳凝聚层的因子化范畴是
\[
\IndCoh^{*}(\Op^{\monfree}_{\check G})\in \FactAlgCat .
\]
它经卷积带有 $\Sph^{\spec}_{\check G}\in \Alg(\FactAlgCat)$ 的作用（该作用的可因子化
构造见 \cite[附录 E.8]{ABC+24a}）。
\end{definition}

\begin{remark}
\label{rem:monfree}
关于无单值 opers 的更详尽讨论，见 \cite[3 节]{ABC+24a}。
\end{remark}

也可沿 $\Ran_{x_0}\to \Ran$ 拉回 $\Op^{\mer}_{\check G}$，得到因子化
$\Op^{\mer}_{\check G}$-模空间
$\Op^{\mer}_{\check G,\Ran_{x_0}}$。
\srcpage{76}

于是
\[
\Op^{\mer}_{\check G,\Ran_{x_0}}
\times_{\LS^{\mer}_{\check G,\Ran_{x_0}}}
\check{\mathfrak{n}}/\check B^{\Ran_{x_0}}
\]
自然带有因子化 $\Op^{\monfree}_{\check G}$-模空间结构。这赋予
\[
\IndCoh^{*}\Bigl(\Op^{\mer}_{\check G,\Ran_{x_0}}
\times_{\LS^{\mer}_{\check G,\Ran_{x_0}}}
\check{\mathfrak{n}}/\check B^{\Ran}_{x_0}\Bigr)
\]
弱因子化 $\IndCoh^{*}(\Op^{\monfree}_{\check G})$-模范畴结构。然而，由于
\[
\check{\mathfrak{n}}/\check B^{\Ran}_{x_0}
\longrightarrow \LS^{\mer}_{\check G,\Ran_{x_0}}
\]
是 ind-proper 的，
\[
\Op^{\mer}_{\check G,\Ran_{x_0}}
\times_{\LS^{\mer}_{\check G,\Ran_{x_0}}}
\check{\mathfrak{n}}/\check B^{\Ran}_{x_0}
\]
相对 $\Ran_{x_0}$ 是 ind-placid ind-概形，因此弱因子化模结构事实上就是因子化模结构
（细节见 \cite[附录 B.13]{ABC+24a}）。

\begin{definition}
\label{def:ifle-pair}
$\FactModCat_{x_0}$ 中的对
\begin{align*}
\bigl(&\IndCoh^{*}(\Op^{\monfree}_{\check G}),\\
&\IndCoh^{*}(\Op^{\mer}_{\check G}\times_{\LS^{\mer}_{\check G}}
\check{\mathfrak{n}}/\check B)\bigr)
\end{align*}
由因子化范畴 $\IndCoh^{*}(\Op^{\monfree}_{\check G})$ 与因子化
$\IndCoh^{*}(\Op^{\monfree}_{\check G})$-模范畴
\begin{align*}
\IndCoh^{*}\bigl(&\Op^{\mer}_{\check G}\times_{\LS^{\mer}_{\check G}}
\check{\mathfrak{n}}/\check B\bigr)\\
&=\IndCoh^{*}\Bigl(\Op^{\mer}_{\check G,\Ran_{x_0}}
\times_{\LS^{\mer}_{\check G,\Ran_{x_0}}}
\check{\mathfrak{n}}/\check B^{\Ran}_{x_0}\Bigr)
\end{align*}
构成。该对带有幺半对 $(\Sph^{\spec}_{\check G},\Aff^{\spec}_{\check G})$ 的作用。
\end{definition}
